% !TEX root = document.tex

\chapter{\label{chap:related-work}Related work}

This thesis lies at the intersection of dependent type elaboration, higher
order unification, and declarative name binding. Each of these topics has been
studied independently. Higher order unification is commonly used during the
elaboration of dependent types, while scope graphs and Statix provide a
declarative approach to name resolution and type system specification.
Higher Order Abstract Syntax provides another way to represent languages with
binders. This chapter places the approach developed in this thesis in relation
to these areas. \par
\section{Higher order unification and dependent type elaboration}
Higher order unification has a long history in logic programming and type
checking. Miller identified the higher order pattern fragment, which restricts
the arguments of unknown functions to distinct bound variables
\autocite{miller_logic_1991}. Within this fragment, unification remains decidable and
solutions have a most general form. The pattern fragment is therefore useful
for systems that need to solve equations involving terms under binders.

Dependent type checking gives rise to such equations during inference and
elaboration. Mazzoli and Abel study the relation between type checking and
unification and show how type checking can be formulated around the generation
and solution of unification constraints \autocite{mazzoli2016typechecking}. Work on
elaboration for dependent type theory also uses unification to reconstruct
information that is omitted from source programs
\autocite{demoura2015elaboration}. Similar problems occur in practical dependent
type systems such as Agda \autocite{norell2007practical}.

Restricting all generated constraints to the pattern fragment is not always
possible. Abel and Pientka describe dynamic pattern unification, where a
constraint that cannot currently be solved may be delayed until more
information becomes available \autocite{abel_extensions_2018}. Their work also
describes pruning, which can remove dependencies that prevent a constraint
from satisfying the pattern restrictions. These techniques provide a more
complete treatment of higher order constraints than the implementation in
this thesis. The Statix implementation presented here handles constraints that
can already be recognized as solvable pattern problems. It does not currently
provide the postponement mechanism required for dynamic pattern unification.

The aim of this thesis is therefore not to introduce a new general higher
order unification procedure. Instead, it studies how the information required
by higher order pattern unification can be obtained when name binding is
represented through scope graphs. In particular, the implementation uses the
scope structure to determine which bound variables may occur in a solution.
It also introduces an explicit representation for the state associated with
unknown terms during inference.

\section{Scope graphs and dependent typing in Statix}

Scope graphs were introduced as a general representation of name resolution
\autocite{neron_extensible_2015}. A scope graph represents scopes and the
visibility relations between them explicitly. This allows the same name
resolution framework to describe both lexical binding and relations introduced
by constructs such as modules and imports. Statix builds on scope graphs to
describe name resolution together with typing constraints
\autocite{van_antwerpen_scopes_2018, 4bf44aa1779c4a968c555e1b54e16119}.

Brouwer previously investigated the implementation of a dependently typed
language in Statix \autocite{none_dependently_2023}. That work shows how core
features of dependent typing can be expressed in the Statix setting. However,
it does not include higher order unification. The present work continues from
this setting and focuses on the inference problems that require higher order
solutions. The main additional problem is that a solution may depend on
variables introduced by surrounding binders. The implementation therefore
needs to relate the context of a unification problem to the binding
information represented in the scope graph.

This also creates a difference from most presentations of higher order
unification. In those presentations, the context of an unknown term is usually
represented directly as part of the unification problem. For example,
contextual representations can associate an open term with the variables on
which it may depend \autocite{nanevski2008contextual}. In Statix, the corresponding
information is already present in the scope graph. This thesis investigates
how that representation can be reused during higher order pattern
unification instead of introducing a separate representation of lexical
binding.

\section{Higher Order Abstract Syntax and logic programming}

Higher Order Abstract Syntax provides a different approach to the
representation of binders. In HOAS, binders in the object language are
represented using binders in the meta language
\autocite{pfenning_higher-order_nodate}. As a consequence, operations involving lexical
binding can be delegated to the underlying language rather than implemented
explicitly.

Logic programming languages based on lambda tree syntax make this approach
executable. Elpi is an implementation of this style of logic programming and
provides support for terms containing binders and for higher order pattern
unification \autocite{dunchev_elpi_2015}. This makes it possible to express typing
rules for lambda abstractions and dependent function types using the binding
and context mechanisms of the meta language.

The Elpi implementation in Chapter~6 is used as a point of comparison for the
Statix implementation. The comparison concerns the representation of binding
rather than the higher order unification algorithm itself. Elpi provides
direct support for lexical binders, substitution, and higher order
unification. Statix represents binding explicitly and requires more of these
operations to be implemented by the language specification. At the same time,
scope graphs directly represent visibility relations that do not follow
lexical nesting, such as module openings and qualified references.

The two approaches therefore emphasize different aspects of language
implementation. The work in this thesis studies higher order inference in a
setting where the explicit name resolution structure provided by scope graphs
is retained, rather than replacing that structure with a higher order
representation of all binders.